%=====================================================================
%  Part 2 opener --- Finding and Synthesising Knowledge
%=====================================================================
\structuralfigures{P2.}

You cannot know whether your question is novel, and you cannot state a
delta, without knowing what has already been established. That is a claim about
a literature --- and like any other claim it needs a method, or it is just an
impression formed from whatever you happened to read. This part treats the
review as a study in its own right: with a protocol, a reproducible search, an
appraisal stage, and a synthesis that produces something that did not exist in
any single paper.

\begin{conceptbox}[What Part~2 Covers]
\begin{tight}
  \item \textbf{Chapter 5 --- Systematic Literature Review.} Protocol, search strings, screening, quality appraisal, extraction, synthesis, PRISMA.
  \item \textbf{Chapter 6 --- Grey Literature and Scoping Reviews.} Snowballing, multivocal reviews, and when the practitioner literature is the evidence.
  \item \textbf{Chapter 7 --- Reference Management and Bibliometrics.} Zotero and BibTeX in practice; mapping a field; why the metrics community distrusts its own metrics.
  \item \textbf{Chapter 8 --- Reading and Appraising a Paper.} The three-pass method, critical appraisal, and reproducing one claim to see whether it holds.
\end{tight}
\end{conceptbox}
